\documentclass[12pt]{article}
\usepackage[a4paper, left=3cm, right=3cm]{geometry}	% Per i margini
\usepackage[latin1]{inputenc}
\usepackage{amssymb}
\usepackage{color, xcolor}
\usepackage{listings}	% Per il codice.
\usepackage{graphicx}	% Per immagini e pdf.
\usepackage{amsthm}		% Per teoremi e dimostrazioni.
\usepackage{amsmath}
\usepackage[backend=bibtex, style=numeric, sorting=none, backref]{biblatex}


% --- Custom colors definitions --- %
\definecolor{mygreen}{rgb}{0,0.6,0}
\definecolor{mygray}{rgb}{0.5,0.5,0.5}


% --- LSTLISTING SETTING --- %
\lstset{ %
	backgroundcolor=\color{white},   			% choose the background color.
	basicstyle=\footnotesize\ttfamily,        % the size of the fonts that are used for the code
	breakatwhitespace=false,         			% sets if automatic breaks should only happen at whitespace
	breaklines=true,                 			% sets automatic line breaking
	captionpos=b,                    			% sets the caption-position to bottom
	commentstyle=\color{mygray},    			% comment style
	deletekeywords={...},            			% if you want to delete keywords from the given language
	escapeinside={\%*}{*)},          			% if you want to add LaTeX within your code
	extendedchars=true,              			% lets you use non-ASCII characters; for 8-bits encodings only, does not work with UTF-8
	frame=single,	                   			% adds a frame around the code
	keepspaces=true,                 			% keeps spaces in text, useful for keeping indentation of code (possibly needs columns=flexible)
	keywordstyle=\color{blue},       		% keyword style
	language=,                 				% the language of the code
	morekeywords={*,...},            			% if you want to add more keywords to the set
	numbers=left,                    			% where to put the line-numbers; possible values are (none, left, right)
	numbersep=5pt,                   			% how far the line-numbers are from the code
	numberstyle=\tiny\color{mygray}, 			% the style that is used for the line-numbers
	rulecolor=\color{black},         			% if not set, the frame-color may be changed on line-breaks within not-black text (e.g. comments (green here))
	showspaces=false,                			% show spaces everywhere adding particular underscores; it overrides 'showstringspaces'
	showstringspaces=false,          			% underline spaces within strings only
	showtabs=false,                  			% show tabs within strings adding particular underscores
	stepnumber=1,                    			% the step between two line-numbers. If it's 1, each line will be numbered
	stringstyle=\color{black},     			% string literal style
	tabsize=2,	                   			% sets default tabsize to 2 spaces
	title=\lstname                   			% show the filename of files included with \lstinputlisting; also try caption instead of title
}

\title{Relazione Sistemi Multi-Agente}
\author{Flavio Bizzarri, Riccardo Grieco, Marco Matarese}

\begin{document}

%\maketitle
% --- TITLE PAGE --- %

\thispagestyle{empty}

\begin{center}
	\includegraphics[width=1\textwidth]{DIETI}
	
	\par\bigskip\par\bigskip\par\bigskip\par\bigskip\par\bigskip\par\bigskip\par\bigskip\par	% Per spazi
	
	{\huge Relazione Progetto di\\}
	{\huge Sistemi Multi-Agente\\}
	
	\par\bigskip\par\bigskip\par\bigskip\par\bigskip\par\bigskip\par\bigskip\par\bigskip\par	% Per spazi
	
	{\LARGE Flavio Bizzarri N97/0000281\\}
	{\LARGE Riccardo Grieco N97/0000286\\}
	{\LARGE Marco Matarese N97/0000280\\}
	
	\par\bigskip\par\bigskip\par\bigskip\par\bigskip\par\bigskip\par\bigskip	% Per spazi
	
	{\large A.A 2017/2018\par}
	
	%\includegraphics[width=0.5\textwidth]{}
\end{center}

%\newpage
\newpage

\section{Architettura della Simulazione}
	\subsection{L'Ambiente}
	\subsection{La Vista}
	\subsection{Il Modello}

\section{Agenti \& Targets}
	\subsection{Comportamento degli Agenti}
	\subsection{Comportamento dei Targets}
	\subsection{L'identificazione dei Targets}
	\subsection{Il Meccanismo d'Asta}
	
\section{Risultati}


\newpage

\tableofcontents


\end{document}
